% ----------------------------------
%   Conclusions and references
% ----------------------------------

\begin{frame}
\frametitle{Conclusions}
\begin{scriptsize}
\begin{itemize}
\item We built the pricing theory from stochastic calculus, through Black-Scholes and arbitrage theory,
up to the Dupire forward equation and its no-arbitrage conditions.
\item Calibration of local volatility is an ill-posed inverse problem,
and Tikhonov regularization with an adjoint gradient solves it at the cost of two PDE solves per iteration.
\item We reproduced the examples of Egger and Engl:
the data determine the volatility near the money, the prior determines the wings.
\item We extended the framework to a time dependent coefficient and calibrated a full SPY surface,
arbitrage-free by construction, with sub-percent reprice errors.
\item The rule of two is confirmed empirically on the calibrated surface,
and the model reprices a held-out maturity within $2\%$.
\end{itemize}
\medskip

\textbf{Further work}
\begin{itemize}
\item Convergence results for the time dependent extension.
\item Inverse problem formulations for random data, closer to market conditions.
\end{itemize}
\end{scriptsize}
\end{frame}

\frame
{\frametitle
{References}
\scriptsize{
\begin{enumerate}
\item Bergomi, L. (2015). \emph{Stochastic Volatility Modeling}. CRC Press.
\item Berestycki, H., Busca, J., Florent, I. (2002). \emph{Asymptotics and calibration of local volatility models}. Quantitative Finance.
\item Bj\"ork, T. (2009). \emph{Arbitrage Theory in Continuous Time}. Oxford University Press.
\item Black, F., Scholes, M. (1973). \emph{The Pricing of Options and Corporate Liabilities}. Journal of Political Economy.
\item Breeden, D. T., Litzenberger, R. H. (1978). \emph{Prices of State-Contingent Claims Implicit in Option Prices}. Journal of Business.
\item Derman, E., Miller, M. B., Park, D. (2016). \emph{The Volatility Smile}. Wiley.
\item Dupire, B. (1994). \emph{Pricing with a Smile}. Risk.
\item Egger, H., Engl, H. W. (2005). \emph{Tikhonov regularization applied to the inverse problem of option pricing: convergence analysis and rates}. Inverse Problems.
\item Engl, H. W., Hanke, M., Neubauer, A. (1996). \emph{Regularization of Inverse Problems}. Kluwer.
\item Merton, R. C. (1973). \emph{Theory of Rational Option Pricing}. Bell Journal of Economics and Management Science.
\item Steele, J. M. (2010). \emph{Stochastic Calculus and Financial Applications}. Springer.
\end{enumerate}
}
}

\begin{frame}
\begin{center}
\Large Thank you
\end{center}
\end{frame}
